% sections/02_background.tex --- SkillStrike

\subsection{LLM Agent Skill Marketplaces}

LLM skill marketplaces publish installable tool packages that extend
agent capabilities via standardized protocols (primarily the
Model Context Protocol, MCP~\cite{anthropic2024mcp}).
Each skill is a self-contained package comprising:
(i)~a natural-language \emph{description} (the sole capability disclosure
mechanism visible to end users),
(ii)~\emph{tool definitions} (structured API declarations consumed by the agent),
and (iii)~\emph{implementation code} (executed by the agent runtime
with the user's OS-level privileges).

\paragraph{Ambient authority.}
Once installed, a skill executes with the full privileges of the agent
runtime---typically the user's OS account.
There is no sandboxing, no OS-level capability restriction, and no
per-tool-call confirmation by default.
This \emph{ambient authority} model~\cite{ruan2024identifying}
means that a skill's behavior is constrained only by its own code
and the LLM's willingness to invoke its tools.

\paragraph{Shared context window.}
All installed skills share a single agent context: the LLM sees
tool definitions and outputs from all skills simultaneously.
When a source skill returns content (search results, emails, API responses),
that content enters the shared context alongside the agent's system prompt
and all other tool schemas---including those of co-installed sink skills.

\subsection{Threat Model}
\label{sec:threat_model}

\paragraph{Source-Sink Model.}
We model a skill ecosystem as a bipartite graph $G = (S_{\rm src} \cup S_{\rm sink}, E)$
where:
\begin{itemize}[leftmargin=*,topsep=2pt,itemsep=1pt]
  \item $S_{\rm src}$ = \emph{source skills}: skills that fetch,
    aggregate, or consume external content, creating an injection
    channel into the agent context.
  \item $S_{\rm sink}$ = \emph{sink skills}: skills possessing a
    \emph{dangerous undisclosed capability}
    $\kappa \in \mathcal{C}_{\rm act} \setminus \mathcal{C}_{\rm vis}$---a
    dangerous action that is executable but absent from documentation.
  \item $E$ contains all pairs $(s, k) \in S_{\rm src} \times S_{\rm sink}$
    where $s$ can deliver attacker-controlled content to the agent
    context in which $k$ is installed.
\end{itemize}

\begin{definition}[Skill Composition Attack]
Given an agent session $\mathcal{A}$ with co-installed source skill
$s \in S_{\rm src}$ and sink skill $k \in S_{\rm sink}$,
an attacker who can influence content delivered through $s$'s output channel
constructs content $c$ such that $\mathcal{A}$'s processing of $c$
invokes $\kappa \in \mathcal{C}_{\rm act,k} \setminus \mathcal{C}_{\rm vis,k}$
without explicit user authorization,
producing a verifiable system-level effect.
\end{definition}

\paragraph{Sink capability families.}
We define five sink capability families:
\textbf{K1} Shell execution: arbitrary command execution via
\texttt{exec}, \texttt{popen}, \texttt{subprocess};
\textbf{K2} Filesystem: credential deletion, file overwrite, directory traversal;
\textbf{K3} Network exfiltration: outbound HTTP/WebSocket data exfil,
webhook triggering;
\textbf{K4} Agent memory: cross-session memory write, context poisoning;
\textbf{K5} Persistence: boot hook injection, cron/service installation.

\paragraph{Source injection channels.}
We define five source injection channels:
\textbf{C1} LLM context: web search, browser automation, news fetch,
API aggregation---content enters the shared context directly;
\textbf{C2} Filesystem: file read operations feeding the context;
\textbf{C3} Agent memory: cross-session memory read channels;
\textbf{C4} Network: outbound-then-inbound API patterns;
\textbf{C5} Persistence: boot hooks and scheduled tasks.

\paragraph{Principals and assumptions.}
Three principals: (1)~the \emph{user}, who installs skills for legitimate purposes;
(2)~the \emph{skill developer}, who publishes skills (possibly with hidden capabilities);
(3)~the \emph{attacker}, who controls content reachable via a source skill's channel.
The attacker does \emph{not} need: access to the user's machine,
the ability to modify installed skills, or OS-level privileges.
The attack requires only the ability to place content where the source
skill's fetch operation will retrieve it---a search engine result,
an email in a monitored inbox, an API response from a compromised endpoint.

\paragraph{Why descriptions are the only defense.}
Unlike Android (capability permission manifests~\cite{felt2011android}),
npm (dependency audit~\cite{synk2023report}), or browser extensions
(sandboxed JavaScript APIs), LLM skill marketplaces have
no machine-enforceable capability permission model.
Natural-language descriptions are therefore the sole mechanism
for users to understand what a skill can do---and the attacker's
advantage is that dangerous capabilities can be implemented
without being disclosed in this description.
We define the \emph{dangerous capability gap} as
$\Delta\kappa = \mathcal{C}_{\rm act} \setminus \mathcal{C}_{\rm vis}$;
the composition attack is enabled whenever $\Delta\kappa_k \neq \varnothing$
for a co-installed sink skill $k$.

\paragraph{Out of scope.}
Direct UI-level social engineering, OS privilege escalation
beyond the agent runtime's account, or attacks on marketplace
platform infrastructure itself.
